% Shared preamble of manuscript.tex and supplement.tex (built with tectonic; see paper/build.py)
\usepackage[a4paper,margin=2.5cm]{geometry}
\usepackage{fontspec}   % XeTeX (tectonic): Latin Modern OpenType by default, Unicode input
\usepackage{microtype}
\usepackage[running]{lineno}
\usepackage{graphicx}
\usepackage{booktabs}
\usepackage{longtable}
\usepackage{tabularx}
\usepackage{array}
\usepackage{caption}
\captionsetup{font=small,labelfont=bf,labelsep=period}
\usepackage{enumitem}
\setlist{itemsep=1pt,topsep=3pt}
\usepackage{amsmath}
\usepackage{xcolor}
\usepackage[numbers,sort&compress,square]{natbib}
\usepackage[hidelinks]{hyperref}
\urlstyle{same}
\setlength{\parskip}{3pt}
\renewcommand{\arraystretch}{1.15}
\newcolumntype{L}[1]{>{\raggedright\arraybackslash}p{#1}}
\newcolumntype{Y}{>{\raggedright\arraybackslash}X}

% Numbers: every number taken from the data is printed with \nm{key}; the values come from
% numbers.tex, which paper/make_numbers.py writes from the repository's files (sources in
% NUMBERS_generated.tsv, text-to-key map in NUMBERS.md). An unknown key stops the build.
\makeatletter
\newcommand{\nmdef}[2]{\expandafter\def\csname nm@#1\endcsname{#2}}
\newcommand{\nm}[1]{\ifcsname nm@#1\endcsname\csname nm@#1\endcsname\else\errmessage{unknown number key: #1}\fi}
\makeatother
% Generated by paper/make_numbers.py. Do not edit. Sources: paper/NUMBERS_generated.tsv
\nmdef{titin_len}{35,991}
\nmdef{pkz_len}{45,212}
\nmdef{uniprotkb_gt_titin}{51}
\nmdef{uniprotkb_gt_pkz}{4}
\nmdef{uniprotkb_pe3}{37}
\nmdef{uniprotkb_pe4}{14}
\nmdef{uniprotkb_reviewed}{0}
\nmdef{uniparc_gt_titin}{477}
\nmdef{uniparc_gt_pkz}{65}
\nmdef{ncbi_gt_titin}{472}
\nmdef{ncbi_gt_pkz}{47}
\nmdef{refseq_gt_titin}{349}
\nmdef{union_seqs}{552}
\nmdef{union_gt_pkz}{71}
\nmdef{union_live}{523}
\nmdef{excl_nolive}{29}
\nmdef{excl_patent}{6}
\nmdef{excl_partial}{7}
\nmdef{longest_record}{187,500}
\nmdef{n_retained}{510}
\nmdef{n_excluded}{42}
\nmdef{n_retained_gt_pkz}{57}
\nmdef{loci_retained_stage0}{228}
\nmdef{xp_records}{329}
\nmdef{xp_genes}{52}
\nmdef{n_euk}{452}
\nmdef{n_bact}{57}
\nmdef{n_arch}{1}
\nmdef{big_euk}{34}
\nmdef{big_bact}{23}
\nmdef{s1_all_A}{114}
\nmdef{s1_all_B}{213}
\nmdef{s1_all_C}{154}
\nmdef{s1_all_D}{29}
\nmdef{s1_big_A}{0}
\nmdef{s1_big_B}{4}
\nmdef{s1_big_C}{36}
\nmdef{s1_big_D}{17}
\nmdef{s2_all_A}{119}
\nmdef{s2_all_B}{267}
\nmdef{s2_all_C}{95}
\nmdef{s2_all_D}{29}
\nmdef{s2_big_A}{0}
\nmdef{s2_big_B}{5}
\nmdef{s2_big_C}{35}
\nmdef{s2_big_D}{17}
\nmdef{fin_all_A}{120}
\nmdef{fin_all_B}{267}
\nmdef{fin_all_C}{94}
\nmdef{fin_all_D}{29}
\nmdef{fin_big_A}{1}
\nmdef{fin_big_B}{5}
\nmdef{fin_big_C}{34}
\nmdef{fin_big_D}{17}
\nmdef{s1s2_C_to_B}{59}
\nmdef{s1s2_B_to_A}{5}
\nmdef{s1s2_C_to_B_apostichopus}{48}
\nmdef{rule_D1_all}{25}
\nmdef{rule_D1_big}{15}
\nmdef{rule_D2_all}{16}
\nmdef{rule_D2_big}{10}
\nmdef{rule_D3_all}{24}
\nmdef{rule_D3_big}{15}
\nmdef{rule_D5_all}{3}
\nmdef{rule_D5_big}{2}
\nmdef{rule_D1_big_euk}{14}
\nmdef{d1_pattern_all}{72}
\nmdef{d1_pattern_big}{32}
\nmdef{d1_pattern_big_euk}{29}
\nmdef{repeat_flag_all}{47}
\nmdef{repeat_flag_big}{17}
\nmdef{flag_big_weak_coding}{11}
\nmdef{flag_big_weak_coding_range}{0.919--0.942}
\nmdef{flag_big_coding}{4}
\nmdef{flag_big_nodna}{2}
\nmdef{d_big_d1_with_dna}{15}
\nmdef{d_big_d5}{2}
\nmdef{rna_not_tested_d1}{0}
\nmdef{rep_models}{43}
\nmdef{rep_testable}{10}
\nmdef{rep_testable_big}{5}
\nmdef{rep_jobs}{7}
\nmdef{rep_runs}{33}
\nmdef{rep_gb}{135}
\nmdef{rep_schisto1_supported}{32}
\nmdef{rep_schisto1_introns}{45}
\nmdef{rep_schisto1_testable}{32}
\nmdef{rep_schisto1_testable_unsup}{0}
\nmdef{rep_schisto1_unique_reads}{44,113}
\nmdef{rep_schisto1_verdict}{partial}
\nmdef{rep_schisto1_grade}{B}
\nmdef{rep_schisto2_supported}{32}
\nmdef{rep_schisto2_introns}{46}
\nmdef{rep_schisto2_testable}{32}
\nmdef{rep_schisto2_testable_unsup}{0}
\nmdef{rep_schisto2_unique_reads}{44,112}
\nmdef{rep_schisto2_verdict}{partial}
\nmdef{rep_schisto2_grade}{B}
\nmdef{rep_owenia_supported}{9}
\nmdef{rep_owenia_introns}{22}
\nmdef{rep_owenia_testable}{6}
\nmdef{rep_owenia_testable_unsup}{0}
\nmdef{rep_owenia_unique_reads}{673}
\nmdef{rep_owenia_verdict}{weak (unsupported introns lie where reads are too few)}
\nmdef{rep_owenia_grade}{C}
\nmdef{rep_arcto_supported}{1}
\nmdef{rep_arcto_introns}{7}
\nmdef{rep_arcto_testable}{0}
\nmdef{rep_arcto_testable_unsup}{0}
\nmdef{rep_arcto_unique_reads}{85}
\nmdef{rep_arcto_verdict}{weak (unsupported introns lie where reads are too few)}
\nmdef{rep_arcto_grade}{C}
\nmdef{rep_noexpr}{6}
\nmdef{rep_noexpr_max_unique}{34}
\nmdef{rep_contradicted}{0}
\nmdef{rep_to_B}{2}
\nmdef{rep_to_D}{0}
\nmdef{rep_run_results}{33}
\nmdef{rep_run_skips}{0}
\nmdef{rep_schisto_err10016938}{32}
\nmdef{mapp_gas1_uj}{25}
\nmdef{mapp_gas1_n}{2,053}
\nmdef{mapp_gas1_pct}{1}
\nmdef{mapp_gas2_uj}{26}
\nmdef{mapp_gas2_n}{1,273}
\nmdef{mapp_gas2_pct}{2}
\nmdef{mapp_gas3_uj}{27}
\nmdef{mapp_gas3_n}{1,208}
\nmdef{mapp_gas3_pct}{2}
\nmdef{mapp_acn_uj}{48}
\nmdef{mapp_acn_n}{1,035}
\nmdef{mapp_acn_pct}{5}
\nmdef{mapp_syn1_uj}{16}
\nmdef{mapp_syn1_n}{203}
\nmdef{mapp_syn1_pct}{8}
\nmdef{mapp_syn2_uj}{9}
\nmdef{mapp_syn2_n}{134}
\nmdef{mapp_syn2_pct}{7}
\nmdef{mapp_schisto1_uj}{37}
\nmdef{mapp_schisto1_n}{45}
\nmdef{mapp_schisto1_pct}{82}
\nmdef{mapp_owenia_uj}{22}
\nmdef{mapp_owenia_n}{22}
\nmdef{mapp_owenia_pct}{100}
\nmdef{mapp_arcto_uj}{7}
\nmdef{mapp_arcto_n}{7}
\nmdef{mapp_arcto_pct}{100}
\nmdef{ctl_release}{2026\_04}
\nmdef{ctl_n}{176}
\nmdef{ctl_silk_n}{3}
\nmdef{ctl_silk_d1}{1}
\nmdef{ctl_spidroin_n}{4}
\nmdef{ctl_spidroin_d1}{3}
\nmdef{ctl_mucin_n}{15}
\nmdef{ctl_mucin_d1}{0}
\nmdef{ctl_adhesin_n}{13}
\nmdef{ctl_adhesin_d1}{4}
\nmdef{ctl_long_n}{150}
\nmdef{ctl_long_d1}{3}
\nmdef{ctl_d1_flagged}{10}
\nmdef{ctl_d23_flagged}{0}
\nmdef{ctl_flagged_with_cds}{8}
\nmdef{ctl_flagged_fickett_min}{1.094}
\nmdef{ctl_flagged_open_max}{1}
\nmdef{hold_n}{9,428}
\nmdef{hold_trigger}{17}
\nmdef{hold_d1_v1}{17}
\nmdef{hold_d1_cds}{10}
\nmdef{hold_d1_v2}{0}
\nmdef{power_models}{57}
\nmdef{power_reps}{20}
\nmdef{power_shift_0.1_v1_contradicted}{0}
\nmdef{power_shift_0.1_v1_under-supported}{0}
\nmdef{power_shift_0.1_v1_partial}{0.997}
\nmdef{power_shift_0.1_v1_full}{0.003}
\nmdef{power_shift_0.1_v2_contradicted}{0}
\nmdef{power_shift_0.1_v2_under-supported}{0}
\nmdef{power_shift_0.1_v2_partial}{0.997}
\nmdef{power_shift_0.1_v2_full}{0.003}
\nmdef{power_shift_0.2_v1_contradicted}{0}
\nmdef{power_shift_0.2_v1_under-supported}{0}
\nmdef{power_shift_0.2_v1_partial}{1}
\nmdef{power_shift_0.2_v1_full}{0}
\nmdef{power_shift_0.2_v2_contradicted}{0}
\nmdef{power_shift_0.2_v2_under-supported}{0.865}
\nmdef{power_shift_0.2_v2_partial}{0.135}
\nmdef{power_shift_0.2_v2_full}{0}
\nmdef{power_shift_0.3_v1_contradicted}{0}
\nmdef{power_shift_0.3_v1_under-supported}{0}
\nmdef{power_shift_0.3_v1_partial}{1}
\nmdef{power_shift_0.3_v1_full}{0}
\nmdef{power_shift_0.3_v2_contradicted}{0}
\nmdef{power_shift_0.3_v2_under-supported}{1}
\nmdef{power_shift_0.3_v2_partial}{0}
\nmdef{power_shift_0.3_v2_full}{0}
\nmdef{power_shift_0.4_v1_contradicted}{0}
\nmdef{power_shift_0.4_v1_under-supported}{0}
\nmdef{power_shift_0.4_v1_partial}{1}
\nmdef{power_shift_0.4_v1_full}{0}
\nmdef{power_shift_0.4_v2_contradicted}{0}
\nmdef{power_shift_0.4_v2_under-supported}{1}
\nmdef{power_shift_0.4_v2_partial}{0}
\nmdef{power_shift_0.4_v2_full}{0}
\nmdef{power_shift_0.5_v1_contradicted}{0.725}
\nmdef{power_shift_0.5_v1_under-supported}{0.198}
\nmdef{power_shift_0.5_v1_partial}{0.076}
\nmdef{power_shift_0.5_v1_full}{0}
\nmdef{power_shift_0.5_v2_contradicted}{0.725}
\nmdef{power_shift_0.5_v2_under-supported}{0.275}
\nmdef{power_shift_0.5_v2_partial}{0}
\nmdef{power_shift_0.5_v2_full}{0}
\nmdef{power_shift_0.6_v1_contradicted}{1}
\nmdef{power_shift_0.6_v1_under-supported}{0}
\nmdef{power_shift_0.6_v1_partial}{0}
\nmdef{power_shift_0.6_v1_full}{0}
\nmdef{power_shift_0.6_v2_contradicted}{1}
\nmdef{power_shift_0.6_v2_under-supported}{0}
\nmdef{power_shift_0.6_v2_partial}{0}
\nmdef{power_shift_0.6_v2_full}{0}
\nmdef{power_skip_0.2_v1_contradicted}{0}
\nmdef{power_skip_0.2_v1_under-supported}{0}
\nmdef{power_skip_0.2_v1_partial}{1}
\nmdef{power_skip_0.2_v1_full}{0}
\nmdef{power_skip_0.2_v2_contradicted}{0}
\nmdef{power_skip_0.2_v2_under-supported}{0.278}
\nmdef{power_skip_0.2_v2_partial}{0.722}
\nmdef{power_skip_0.2_v2_full}{0}
\nmdef{power_skip_0.3_v1_contradicted}{0}
\nmdef{power_skip_0.3_v1_under-supported}{0}
\nmdef{power_skip_0.3_v1_partial}{1}
\nmdef{power_skip_0.3_v1_full}{0}
\nmdef{power_skip_0.3_v2_contradicted}{0}
\nmdef{power_skip_0.3_v2_under-supported}{1}
\nmdef{power_skip_0.3_v2_partial}{0}
\nmdef{power_skip_0.3_v2_full}{0}
\nmdef{power_skip_0.5_v1_contradicted}{0.242}
\nmdef{power_skip_0.5_v1_under-supported}{0.061}
\nmdef{power_skip_0.5_v1_partial}{0.696}
\nmdef{power_skip_0.5_v1_full}{0}
\nmdef{power_skip_0.5_v2_contradicted}{0.242}
\nmdef{power_skip_0.5_v2_under-supported}{0.758}
\nmdef{power_skip_0.5_v2_partial}{0}
\nmdef{power_skip_0.5_v2_full}{0}
\nmdef{power_skip_0.6_v1_contradicted}{0.981}
\nmdef{power_skip_0.6_v1_under-supported}{0.009}
\nmdef{power_skip_0.6_v1_partial}{0.011}
\nmdef{power_skip_0.6_v1_full}{0}
\nmdef{power_skip_0.6_v2_contradicted}{0.981}
\nmdef{power_skip_0.6_v2_under-supported}{0.019}
\nmdef{power_skip_0.6_v2_partial}{0}
\nmdef{power_skip_0.6_v2_full}{0}
\nmdef{pw_contr_ge60_min_pct}{98}
\nmdef{pw_shift50_pct}{73}
\nmdef{pw_skip50_pct}{24}
\nmdef{pw_v1_partial_2040_min_pct}{99.9}
\nmdef{pw_v2_down_3040_min_pct}{100}
\nmdef{pw_v2_shift20_pct}{87}
\nmdef{pw_v2_skip20_pct}{28}
\nmdef{pw_v2_10_down_pct}{0}
\nmdef{dnactl_v3_a_deletion_ill}{against: 1 single-copy boundary(ies) with 0 spanning pairs between well-covered windows; 1 clip pile-up(s) where physical coverage drops; discordant pairs up to 19\% in a window [0.66 kb in 1 stretch(es) where this library cannot place pairs uniquely (repeat), not crossed]}
\nmdef{dnactl_v3_a_deletion_illont}{against: 1 single-copy boundary(ies) with 0 spanning pairs between well-covered windows; 1 clip pile-up(s) where physical coverage drops; discordant pairs up to 19\% in a window [0.66 kb in 1 stretch(es) where this library cannot place pairs uniquely (repeat), not crossed]}
\nmdef{dnactl_v3_b_collapse_ill}{supported by the reads}
\nmdef{dnactl_v3_b_collapse_illont}{supported by the reads}
\nmdef{dnactl_v3_c_expansion_ill}{supported outside 9.37 kb of repeat (no long reads to resolve it)}
\nmdef{dnactl_v3_c_expansion_illont}{supported outside 9.37 kb of repeat; long reads do not span it}
\nmdef{dnactl_v3_d_rrna_correct_ill}{supported outside 4.89 kb of repeat (no long reads to resolve it)}
\nmdef{dnactl_v3_d_rrna_correct_illont}{supported by the reads (short reads; the 4.89 kb repeat block(s) spanned by 13+ long reads)}
\nmdef{dnactl_v3_e_segdup_correct_ill}{inconclusive: only 98.7\% of single-copy gene bases at depth $\geq$ 5 [12.0 kb in multi-copy block(s) that short reads cannot place; 13.04 kb in 4 stretch(es) where this library cannot place pairs uniquely (repeat), not crossed]}
\nmdef{dnactl_v3_e_segdup_correct_illont}{inconclusive: only 98.7\% of single-copy gene bases at depth $\geq$ 5 [12.0 kb in multi-copy block(s) that short reads cannot place; 13.04 kb in 4 stretch(es) where this library cannot place pairs uniquely (repeat), not crossed]}
\nmdef{dnactl_v31_a_deletion_ill}{against: 1 single-copy boundary(ies) with 0 spanning pairs between well-covered windows; 1 clip pile-up(s) where physical coverage drops; discordant pairs up to 19\% in a window [0.66 kb in 1 stretch(es) where this library cannot place pairs uniquely (repeat), not crossed]}
\nmdef{dnactl_v31_a_deletion_illont}{against: 1 single-copy boundary(ies) with 0 spanning pairs between well-covered windows; 1 clip pile-up(s) where physical coverage drops; discordant pairs up to 19\% in a window [0.66 kb in 1 stretch(es) where this library cannot place pairs uniquely (repeat), not crossed]}
\nmdef{dnactl_v31_b_collapse_ill}{against: copy-number anomaly (v3.1): 5 kb at 1.98x, excess copies (collapsed repeat?)}
\nmdef{dnactl_v31_b_collapse_ill_ratio}{1.98}
\nmdef{dnactl_v31_b_collapse_illont}{against: copy-number anomaly (v3.1): 5 kb at 1.98x, excess copies (collapsed repeat?)}
\nmdef{dnactl_v31_b_collapse_illont_ratio}{1.98}
\nmdef{dnactl_v31_c_expansion_ill}{against: copy-number anomaly (v3.1): 10 kb at 0.49x, deficit (expanded or absent?) [8.0 kb in multi-copy block(s) that short reads cannot place; 8.74 kb in 2 stretch(es) where this library cannot place pairs uniquely (repeat), not crossed]}
\nmdef{dnactl_v31_c_expansion_ill_ratio}{0.49}
\nmdef{dnactl_v31_c_expansion_illont}{against: copy-number anomaly (v3.1): 10 kb at 0.49x, deficit (expanded or absent?) [8.0 kb in multi-copy block(s) that short reads cannot place; 8.74 kb in 2 stretch(es) where this library cannot place pairs uniquely (repeat), not crossed]}
\nmdef{dnactl_v31_c_expansion_illont_ratio}{0.49}
\nmdef{dnactl_v31_d_rrna_correct_ill}{supported outside 4.89 kb of repeat (no long reads to resolve it)}
\nmdef{dnactl_v31_d_rrna_correct_illont}{supported by the reads (short reads; the 4.89 kb repeat block(s) spanned by 13+ long reads)}
\nmdef{dnactl_v31_e_segdup_correct_ill}{inconclusive: only 98.7\% of single-copy gene bases at depth $\geq$ 5 [12.0 kb in multi-copy block(s) that short reads cannot place; 13.04 kb in 4 stretch(es) where this library cannot place pairs uniquely (repeat), not crossed]}
\nmdef{dnactl_v31_e_segdup_correct_illont}{inconclusive: only 98.7\% of single-copy gene bases at depth $\geq$ 5 [12.0 kb in multi-copy block(s) that short reads cannot place; 13.04 kb in 4 stretch(es) where this library cannot place pairs uniquely (repeat), not crossed]}
\nmdef{dnactl_real_unchanged}{yes}
\nmdef{dnactl_repeat_depth_ratio_range}{0.84--0.97}
\nmdef{dnactl_omni_repeat_depth_ratio_range}{0.86--0.97}
\nmdef{sens_variants}{33}
\nmdef{sens_big_ab_stable}{31}
\nmdef{sens_big_ab_changed}{D3 open frames $\geq$ 1 (published 2), spliced reads per intron $\geq$ 10 (published 3)}
\nmdef{sens_s6_lenient_all_B}{267}
\nmdef{sens_base_all_B}{267}
\nmdef{sens_d1v1_big_D}{34}
\nmdef{sens_d1v1_all_D}{76}
\nmdef{sens_allA_range}{116--178}
\nmdef{dna_sens_settings}{7}
\nmdef{dna_sens_grade_changes}{min\_pairs\_per\_boundary=10}
\nmdef{prot_spectra}{594,284}
\nmdef{prot_db_entries}{151,457}
\nmdef{prot_fractions}{60}
\nmdef{prot_scan_share_pct}{48}
\nmdef{prot_id_rate_pct}{84.6}
\nmdef{prot_distal}{8,279, 12,982, 15,831, 16,949}
\nmdef{prot_psms}{502,718}
\nmdef{prot_decoy_psms}{5,026}
\nmdef{prot_peptides}{42,000}
\nmdef{prot_proteins}{7,981}
\nmdef{prot_giant_len}{39,686}
\nmdef{prot_giant_peptides}{186}
\nmdef{prot_giant_shared}{0}
\nmdef{prot_giant_psms}{1,487}
\nmdef{prot_giant_residues}{2,357}
\nmdef{prot_giant_cov_pct}{5.9}
\nmdef{prot_giant_regions}{86}
\nmdef{prot_giant_first}{292}
\nmdef{prot_giant_last_end}{16,957}
\nmdef{prot_giant_last_pct}{42.7}
\nmdef{prot_giant_fifths}{182, 3, 1, 0, 0}
\nmdef{prot_giant_first6000}{182}
\nmdef{prot_giant_after6000}{8,279, 12,982, 15,831, 16,949}
\nmdef{prot_giant_protein_q}{0.0002}
\nmdef{prot_rev_q}{0.063}
\nmdef{prot_shuf_q}{0.17}
\nmdef{prot_rev_psms}{0}
\nmdef{prot_shuf_psms}{0}
\nmdef{prot_tryptic_range}{41--63}
\nmdef{prot_repeat_span}{17,653--25,339}
\nmdef{prot_repeat_blocks}{5}
\nmdef{prot_peptides_in_or_after_repeats}{0}
\nmdef{prot_eftu_peptides}{23}
\nmdef{prot_eftu_cov_pct}{55}
\nmdef{prot_eftu_len}{400}
\nmdef{prot_l7_peptides}{6}
\nmdef{prot_l7_cov_pct}{68}
\nmdef{prot_l7_len}{128}
\nmdef{prot_top_n}{12}
\nmdef{prot_top_rank_giant}{11}
\nmdef{prot_top_cov_range}{40.8--78.3}
\nmdef{prot_wall_min}{35}
\nmdef{prot_rss_gb}{5.2}
\nmdef{raw_files}{60}
\nmdef{raw_conversion_errors}{0}
\nmdef{raw_ms2}{1,106,517}
\nmdef{raw_mgf_share_pct}{54}
\nmdef{raw_instrument}{Orbitrap Velos Pro}
\nmdef{raw_psms}{507,612}
\nmdef{raw_peptides}{40,845}
\nmdef{raw_proteins}{7,560}
\nmdef{raw_giant_peptides}{182}
\nmdef{raw_giant_psms}{1,505}
\nmdef{raw_giant_residues}{2,282}
\nmdef{raw_giant_cov_pct}{5.8}
\nmdef{raw_giant_fifths}{179, 3, 0, 0, 0}
\nmdef{raw_giant_first6000}{179}
\nmdef{raw_giant_after6000}{8,279, 12,982, 15,831}
\nmdef{raw_giant_last_end}{15,840}
\nmdef{raw_giant_last_pct}{39.9}
\nmdef{raw_giant_protein_q}{0.00022}
\nmdef{raw_rev_q}{0.089}
\nmdef{raw_shuf_q}{0.17}
\nmdef{raw_rev_psms}{0}
\nmdef{raw_shuf_psms}{0}
\nmdef{raw_eftu_peptides}{24}
\nmdef{raw_eftu_cov_pct}{57}
\nmdef{raw_l7_peptides}{5}
\nmdef{raw_l7_cov_pct}{56}
\nmdef{raw_both}{179}
\nmdef{raw_only_raw}{3}
\nmdef{raw_only_mgf}{7}
\nmdef{raw_mgf_only_after6000}{16,949}
\nmdef{raw_in_or_after_repeats}{0}
\nmdef{raw_wall_min}{56}
\nmdef{prot_peak_last_q}{0.0096}
\nmdef{sy_ms2}{131,021}
\nmdef{sy_raw_files}{5}
\nmdef{sy_db_kept}{99,767}
\nmdef{sy_db_authors}{99,773}
\nmdef{sy_fragments}{6}
\nmdef{sy_fragment_max_len}{46,359}
\nmdef{sy_psms_spec}{272}
\nmdef{sy_peptides_spec}{88}
\nmdef{sy_min_pep_q}{0.018}
\nmdef{sy_min_prot_q}{0.024}
\nmdef{sy_authors_n}{37}
\nmdef{sy_authors_found}{31}
\nmdef{sy_giant_psms_spec}{0}
\nmdef{sy_giant_psms_any}{17}
\nmdef{sy_rev_psms_any}{27}
\nmdef{sy_shuf_psms_any}{15}
\nmdef{sy_calc3_psms}{30}
\nmdef{sy_tub_psms}{37}
\nmdef{sy_ef2_psms}{9}
\nmdef{prot_authors_file}{Mudpit\_160213\_OV3\_P5\_SK\_JHa\_372-Pe1\_01.mzid.gz}
\nmdef{prot_authors_len}{39,671}
\nmdef{prot_authors_peptides}{174}
\nmdef{prot_authors_fifths}{130, 2, 2, 2, 0}
\nmdef{prot_authors_last_start}{25,870}
\nmdef{prot_authors_pass_peptides}{130}
\nmdef{prot_authors_pass_fifths}{101, 0, 0, 0, 0}
\nmdef{pride_omni}{0}
\nmdef{pride_noc881}{0}
\nmdef{pride_noc1730}{0}
\nmdef{pride_sycon}{2}
\nmdef{pride_velam}{1}
\nmdef{prot_raw_files}{60}
\nmdef{prot_raw_gb}{15.3}
\nmdef{prot_closed_genome_proteins}{1,863}
\nmdef{raw_giant_last_pct_round}{40}
\nmdef{rep_untestable_noruns}{28}
\nmdef{rep_untestable_nolocus}{5}
\nmdef{twin_fix_n}{6}
\nmdef{twin_fix_big}{0}
\nmdef{rep_noruns_plectranthus}{25}
\nmdef{rep_noruns_merlangius}{2}
\nmdef{rep_noruns_coilia}{1}
\nmdef{raw_giant_unc_pct}{60}
\nmdef{prot_giant_grade}{C}
\nmdef{lb_loci}{213}
\nmdef{lb_excluded_rows}{42}
\nmdef{lb_A}{9}
\nmdef{lb_B}{85}
\nmdef{lb_C}{90}
\nmdef{lb_D}{29}
\nmdef{lb_big}{55}
\nmdef{lb_big_A}{1}
\nmdef{lb_big_B}{3}
\nmdef{lb_big_C}{34}
\nmdef{lb_big_D}{17}
\nmdef{lb_merged_records}{11}
\nmdef{lb_A_titin_loci}{8}
\nmdef{lb_A_titin_min}{36,247}
\nmdef{lb_A_titin_max}{38,490}
\nmdef{rna_runs_total}{269}
\nmdef{rna_runs_kept}{261}
\nmdef{rna_runs_superseded}{8}
\nmdef{rna_runs_followup}{16}
\nmdef{rna_species}{71}
\nmdef{rna_runs_v1}{8}
\nmdef{rna_runs_v2}{253}
\nmdef{rna_reads_billion}{11.6}
\nmdef{rna_runs_stream_error}{51}
\nmdef{rna_runs_stream_error_first}{46}
\nmdef{rna_stream_error_frac_median_pct}{68}
\nmdef{rna_stream_error_frac_min_pct}{9}
\nmdef{rna_runs_saturation}{55}
\nmdef{rna_runs_readcap}{20}
\nmdef{rna_frac_read_median_pct}{88}
\nmdef{rna_runs_lt50}{49}
\nmdef{rna_full}{59}
\nmdef{rna_partial}{75}
\nmdef{rna_weak}{27}
\nmdef{rna_noexpr}{12}
\nmdef{rna_noruns}{127}
\nmdef{rna_nottested}{15}
\nmdef{rna_notselected}{137}
\nmdef{rna_prok}{58}
\nmdef{rna_contradicted}{0}
\nmdef{rna_tested_models}{173}
\nmdef{titin_runs}{6}
\nmdef{titin_reads_unique}{361,467}
\nmdef{titin_n_introns}{361}
\nmdef{titin_introns_unique_support}{340}
\nmdef{titin_introns_multi_only}{13}
\nmdef{titin_introns_counted_supported}{353}
\nmdef{titin_cds_breadth}{0.973}
\nmdef{titin_testable_introns}{325}
\nmdef{titin_testable_introns_unsupported}{0}
\nmdef{pkz_runs}{7}
\nmdef{pkz_reads_unique}{71,183}
\nmdef{pkz_n_introns}{16}
\nmdef{pkz_introns_unique_support}{9}
\nmdef{pkz_introns_multi_only}{7}
\nmdef{pkz_introns_counted_supported}{16}
\nmdef{pkz_cds_breadth}{0.733}
\nmdef{pkz_testable_introns}{12}
\nmdef{pkz_testable_introns_unsupported}{0}
\nmdef{sycon_runs}{6}
\nmdef{sycon_reads_unique}{49,375}
\nmdef{sycon_n_introns}{24}
\nmdef{sycon_introns_unique_support}{22}
\nmdef{sycon_introns_multi_only}{0}
\nmdef{sycon_introns_counted_supported}{22}
\nmdef{sycon_cds_breadth}{1.0}
\nmdef{sycon_testable_introns}{24}
\nmdef{sycon_testable_introns_unsupported}{2}
\nmdef{halichondria_runs}{4}
\nmdef{halichondria_reads_unique}{70,304}
\nmdef{halichondria_n_introns}{2}
\nmdef{halichondria_introns_unique_support}{2}
\nmdef{halichondria_introns_multi_only}{0}
\nmdef{halichondria_introns_counted_supported}{2}
\nmdef{halichondria_cds_breadth}{0.979}
\nmdef{halichondria_testable_introns}{0}
\nmdef{halichondria_testable_introns_unsupported}{0}
\nmdef{titin_introns_none}{8}
\nmdef{microcaecilia_testable}{41}
\nmdef{microcaecilia_testable_unsup}{1}
\nmdef{microcaecilia_verdict}{partial}
\nmdef{rana_testable}{63}
\nmdef{rana_testable_unsup}{1}
\nmdef{rana_verdict}{partial}
\nmdef{solea_testable}{57}
\nmdef{solea_testable_unsup}{1}
\nmdef{solea_verdict}{partial}
\nmdef{pantherophis_testable}{21}
\nmdef{pantherophis_testable_unsup}{0}
\nmdef{pantherophis_verdict}{weak (unsupported introns lie where reads are too few)}
\nmdef{liftover_moved}{143}
\nmdef{liftover_n_introns}{334}
\nmdef{cam_after_sup}{67}
\nmdef{cam_after_testable_unsup}{0}
\nmdef{ovis_length}{98,182}
\nmdef{ovis_n_exons}{37}
\nmdef{ovis_genomic_span_bp}{356,419}
\nmdef{ovis_protein_period_aa}{22}
\nmdef{ovis_protein_period_identity}{0.847}
\nmdef{ovis_rep25_frac}{0.9922}
\nmdef{ovis_gc_cds}{0.678}
\nmdef{ovis_dna_tandem_lag_bp}{22}
\nmdef{ovis_dna_tandem_identity}{0.939}
\nmdef{ovis_fickett_median}{0.465}
\nmdef{ovis_n_other_frames_open}{2}
\nmdef{ovis_min_stops_per_kcodon_other_frames}{1.2}
\nmdef{ovis_rna_reads}{50,777,537}
\nmdef{ovis_rna_candidates}{0}
\nmdef{ovis_period22_entries}{10}
\nmdef{cirrhinus_len}{59,714}
\nmdef{cirrhinus_exons}{454}
\nmdef{cirrhinus_runs}{6}
\nmdef{cirrhinus_reads_unique}{54,161}
\nmdef{cirrhinus_n_introns}{453}
\nmdef{cirrhinus_introns_counted_supported}{356}
\nmdef{bagarius_len}{54,498}
\nmdef{bagarius_exons}{417}
\nmdef{bagarius_runs}{2}
\nmdef{bagarius_reads_unique}{4,277}
\nmdef{bagarius_n_introns}{416}
\nmdef{bagarius_introns_counted_supported}{144}
\nmdef{cirrhinus_second_copy_start_aa}{31,001}
\nmdef{cirrhinus_first_ttna}{30 of 30}
\nmdef{cirrhinus_second_ttnb}{20 of 29}
\nmdef{bagarius_second_copy_start_aa}{32,001}
\nmdef{bagarius_first_ttna}{31 of 31}
\nmdef{bagarius_second_ttnb}{11 of 22}
\nmdef{cirrhinus_join_intron_aa}{30,675}
\nmdef{cirrhinus_join_intron_reads}{0}
\nmdef{cirrhinus_join_before_median}{103}
\nmdef{cirrhinus_join_after_median}{2}
\nmdef{ovis_entries}{11}
\nmdef{ovis_period_multiple22}{11}
\nmdef{sycon_gap_patch_bp}{2,968}
\nmdef{sycon_best_homolog_len}{4,063}
\nmdef{sycon_len}{50,516}
\nmdef{nbc00881_len}{59,275}
\nmdef{nbc00881_ratio}{1.27}
\nmdef{nbc00881_best_len}{46,734}
\nmdef{nbc00881_orf_bp}{177,828}
\nmdef{nbc01730_len}{51,778}
\nmdef{nbc01730_ratio}{1.11}
\nmdef{nbc01730_best_len}{46,734}
\nmdef{nbc01730_orf_bp}{155,337}
\nmdef{pkz_uniprot_pe}{Predicted}
\nmdef{stage1_cpu_hours}{4.4}
\nmdef{ovis_stops_six_frames}{0, 1.2, 1.2, 45.6, 45.2, 45.5}
\nmdef{dna_oub_reads}{83,321,080}
\nmdef{dna_oub_reads_expected}{83,321,080}
\nmdef{dna_oub_depth}{45}
\nmdef{dna_oub_ratio}{1.12}
\nmdef{dna_oub_span_median}{38}
\nmdef{dna_oub_span_min_testable}{4}
\nmdef{dna_oub_discordant_pct}{1.5}
\nmdef{dna_oub_insert}{288}
\nmdef{dna_oub_untestable}{1}
\nmdef{dna_oub_boundaries}{257}
\nmdef{dna_oub_verdict}{inconclusive: weakest single-copy boundary spanned by 4 [1.34 kb in 2 stretch(es) where this library cannot place pairs uniquely (repeat), not crossed]}
\nmdef{dna_oub_control}{yes}
\nmdef{dna_oub_v2}{against: 1 single-copy boundary(ies) with 0 spanning pairs between well-covered windows}
\nmdef{dna_oub_E_at_untestable}{0}
\nmdef{dna_oub_obs_at_untestable}{0}
\nmdef{dna_oub_weakest}{4--5.8}
\nmdef{dna_nbc00881_reads}{14,056,286}
\nmdef{dna_nbc00881_reads_expected}{14,056,286}
\nmdef{dna_nbc00881_depth}{252}
\nmdef{dna_nbc00881_ratio}{1.02}
\nmdef{dna_nbc00881_span_median}{252}
\nmdef{dna_nbc00881_span_min_testable}{159}
\nmdef{dna_nbc00881_discordant_pct}{1.2}
\nmdef{dna_nbc00881_insert}{340}
\nmdef{dna_nbc00881_untestable}{1}
\nmdef{dna_nbc00881_boundaries}{177}
\nmdef{dna_nbc00881_verdict}{supported by the reads (short reads; the 1.1 kb repeat block(s) spanned by 64+ long reads)}
\nmdef{dna_nbc00881_control}{yes}
\nmdef{dna_nbc00881_v2}{supported by the reads}
\nmdef{dna_nbc00881_E_at_untestable}{0}
\nmdef{dna_nbc00881_obs_at_untestable}{6}
\nmdef{dna_nbc00881_weakest}{159--163}
\nmdef{dna_nbc00881_ont_reads}{61,843}
\nmdef{dna_nbc00881_ont_span_blocks}{64--67}
\nmdef{dna_nbc01730_reads}{15,565,314}
\nmdef{dna_nbc01730_reads_expected}{15,565,314}
\nmdef{dna_nbc01730_depth}{272}
\nmdef{dna_nbc01730_ratio}{1.02}
\nmdef{dna_nbc01730_span_median}{284}
\nmdef{dna_nbc01730_span_min_testable}{74}
\nmdef{dna_nbc01730_discordant_pct}{1.4}
\nmdef{dna_nbc01730_insert}{347}
\nmdef{dna_nbc01730_untestable}{8}
\nmdef{dna_nbc01730_boundaries}{155}
\nmdef{dna_nbc01730_verdict}{supported by the reads (short reads; the 9.0 kb repeat block(s) spanned by 9+ long reads)}
\nmdef{dna_nbc01730_control}{yes}
\nmdef{dna_nbc01730_v2}{supported by the reads (short reads; the 9.0 kb repeat block(s) spanned by 9+ long reads)}
\nmdef{dna_nbc01730_E_at_untestable}{0, 0, 0, 0, 0, 0, 0, 0}
\nmdef{dna_nbc01730_obs_at_untestable}{0, 0, 0, 0, 0, 0, 0, 0}
\nmdef{dna_nbc01730_weakest}{74--22.2}
\nmdef{dna_nbc01730_ont_reads}{36,579}
\nmdef{dna_nbc01730_ont_span_blocks}{9}
\nmdef{dna_tdt_reads}{85,214,796}
\nmdef{dna_tdt_reads_expected}{85,214,796}
\nmdef{dna_tdt_depth}{40}
\nmdef{dna_tdt_ratio}{1.08}
\nmdef{dna_tdt_span_median}{34}
\nmdef{dna_tdt_span_min_testable}{14}
\nmdef{dna_tdt_discordant_pct}{2.1}
\nmdef{dna_tdt_insert}{295}
\nmdef{dna_tdt_untestable}{1}
\nmdef{dna_tdt_boundaries}{257}
\nmdef{dna_tdt_verdict}{supported outside 1.32 kb of repeat (no long reads to resolve it)}
\nmdef{dna_tdt_control}{yes}
\nmdef{dna_tdt_v2}{n/a (judged under v3 only)}
\nmdef{dna_tdt_E_at_untestable}{0}
\nmdef{dna_tdt_obs_at_untestable}{0}
\nmdef{dna_tdt_weakest}{14--19.1}
\nmdef{dna_pmg_reads}{78,636,086}
\nmdef{dna_pmg_reads_expected}{78,636,086}
\nmdef{dna_pmg_depth}{52}
\nmdef{dna_pmg_ratio}{1.06}
\nmdef{dna_pmg_span_median}{45}
\nmdef{dna_pmg_span_min_testable}{15}
\nmdef{dna_pmg_discordant_pct}{1.9}
\nmdef{dna_pmg_insert}{292}
\nmdef{dna_pmg_untestable}{1}
\nmdef{dna_pmg_boundaries}{257}
\nmdef{dna_pmg_verdict}{supported outside 1.3 kb of repeat (no long reads to resolve it)}
\nmdef{dna_pmg_control}{yes}
\nmdef{dna_pmg_v2}{against: 1 single-copy boundary(ies) with 0 spanning pairs between well-covered windows}
\nmdef{dna_pmg_E_at_untestable}{0}
\nmdef{dna_pmg_obs_at_untestable}{0}
\nmdef{dna_pmg_weakest}{15--43.6}
\nmdef{orf_bp}{257,412}
\nmdef{contig_len_tdt}{1,614,705}
\nmdef{rep_lag_bp}{10,362}
\nmdef{rep_lag_codons}{3,454}
\nmdef{rep_copy_bp}{1,184}
\nmdef{rep_copy_a}{223,454--224,637}
\nmdef{rep_copy_b}{233,816--234,999}
\nmdef{rep_diffs}{2}
\nmdef{orf_len_aa}{85,803}
\nmdef{rep_copy_bp_kb}{1.18}
\nmdef{dna_oub_weakest_obs}{4}
\nmdef{dna_oub_weakest_exp}{6}
\nmdef{orf_len_if_excised}{82,349}
\nmdef{omni_copies}{3}
\nmdef{omni_aa_identity_min_pct}{99.96}
\nmdef{omni_aa_identity_max_pct}{99.98}
\nmdef{omni_contig_lengths}{1,614,705, 1,613,271, 1,168,148}
\nmdef{omni_n_in_cds}{0}
\nmdef{hmm_panel}{74}
\nmdef{s2_A_titin}{119}
\nmdef{s1_fix_B_to_C}{48}
\nmdef{titin_downgrades_withdrawn}{5}
\nmdef{cirrhinus_first_end}{31,000}


\newcommand{\gene}[1]{\textit{#1}}
\newcommand{\spn}[1]{\textit{#1}}
\newcommand{\Ca}{\textit{Candidatus}}
\newcommand{\placeholder}[1]{\textcolor{red!70!black}{[#1]}}
% table bodies from paper/tables are read with the primitive \input (LaTeX's \input breaks \bottomrule)
\makeatletter\newcommand{\tablebody}[1]{\@@input #1 }\makeatother
\newcommand{\file}[1]{\texttt{\detokenize{#1}}}
